\section{An explicit finite-arithmetic rate}
\label{sec:closed-rate}

This section gives a root-free expression for the rate on a finite
rational input and describes its bit complexity. Let
\(\lfloor z\rfloor_2\) be the largest integral power of two at most
\(z>0\), and let \(\lfloor0\rfloor_2=0\). The empty-alternative branch
returns zero when \(e\ge1-\min_iq_i\). The minimum is taken on the
declared alphabet, including zero-probability labels.

Assume the alternative is nonempty, set \(\delta=1-e\), and compute
\[
  \widehat J=\sum_i q_i\min\{1,\lfloor\sqrt{q_i/\delta}\rfloor_2\}.
\]
Then \(J/2\le\widehat J\le J\). If \(\widehat J\ge16\delta\), return
\(\lceil1/\widehat J\rceil\). This is the near-maximum-distance
rate, since in this branch \(J\ge16\delta\). In the other branch
\(J<32\delta\), which is sufficient for all the regular-regime
baseline comparisons.

Use the shortest decreasing-mass prefix \(H\) with \(q(H)\ge\delta\).
Its complement \(T\) has mass \(w\). Put \(A=1-w\), \(\Delta=e-w\),
and \(O=H\setminus\{i_*\}\), with \(i_*\) a largest head label.
The head cuts \(R\) are prefixes of \(O\), and the tail cuts \(S\) are
prefixes of \(T\), each in decreasing reference mass. It suffices
to use endpoints of equal-probability runs, including the empty
and full cut in each list. For each \(R\) set
\[
 \begin{split}
 c_{R,0}&=q(O\setminus R),\\
 c_{R,1}&=\max\left\{\lfloor\sqrt{wA}\rfloor_2,\
       \left\lfloor
          \left(\sum_{i\in R}\lfloor q_i^{2/3}\rfloor_2\right)^{3/4}
       \right\rfloor_2\right\},\\
 c_{R,2}&=1.
 \end{split}
\]
For each \(S\) set
\[
 b_S=\max\{\Delta,q(S)\},\qquad
 h_S=\left(\sum_{i\in T\setminus S}q_i\lfloor\sqrt{q_i}\rfloor_2\right)^2 .
\]
A tail consisting solely of null-zero atoms contributes only the
empty cut, with \(b_S=\Delta\) and \(h_S=0\).

For \(c,b\ge0\) and a nonnegative integer \(j\), define
\[
 E(c,b,j)=
 \begin{cases}
 0,&c\le b,\\
 +\infty,&c>b\ \text{and}\ (b=0\ \text{or}\ j=0),\\
 \displaystyle\left\lceil\frac{\lceil\log_2(c/b)\rceil}{j}\right\rceil,
   &\text{otherwise}.
 \end{cases}
\]
The resulting rate in the regular regime is
\begin{equation}
 4^{\,\displaystyle\min_R\,
       \max_{\substack{S\\ a\in\{0,1,2\}}}
       \min\{E(c_{R,a},b_S,a),\,E(c_{R,a},h_S,a+4)\}} .
 \label{eq:closed-rate}
\end{equation}
The use of \(a\) as a numerical index inside \eqref{eq:closed-rate}
is local to that display.

\subsection{Comparison with the statistical crossing}

For any \(n\ge1\) define
\[
 \widehat D(n)=\min_R\max_{a=0,1,2}c_{R,a}n^{-a/2},
 \qquad
 \widehat Y(n)=\min_S\max\{b_S,h_Sn^2\}.
\]
The original head cost for each \(R\) is the sum of four terms:
the omitted mass, the group standard deviation divided by \(\sqrt n\),
the retained bulk divided by \(\sqrt n\), and \(1/n\).
The group numerator loses at most a factor 2 in rounding. The
head numerator loses at most \(2^{7/4}<4\), including the inner
rounding and the final rounding. Thus at every cut the rounded
maximum lies between one eighth of the original cost and the
original cost. Within an equal-probability run, the original head
cost is affine plus a positive multiple of a \(3/4\)-power of an
affine function of the number of retained labels. Concavity places
its minimum at an endpoint. Consequently
\begin{equation}
  D(n)/8\le\widehat D(n)\le D(n).
  \label{eq:closed-head-comparison}
\end{equation}
The rounded maximum itself may prefer a cut within a tied run.
The same constant-factor estimate holds with endpoint compression.
An expanded scan of all tied labels can return a different integer.

For the tail, let \(z=\sqrt{X(n)}\), and put
\(Q_S=q(S)\), \(H_S=\sum_{i\in T\setminus S}q_i^{3/2}\).
If \(z>0\), its defining equation gives
\[
 z^2=\sum_{i\in T}\min(q_i,nq_i^{3/2}z)
       \le Q_S+nH_Sz
       \qquad\text{for every }S.
\]
If \(B_S=\max\{Q_S,(nH_S)^2\}\), this implies
\(z^2\le B_S+\sqrt{B_S}\,z\), hence \(z^2\le4B_S\).
Conversely, equality holds in the previous display for the prefix
\(\{i:nz\sqrt{q_i}\ge1\}\), which then has \(B_S\le z^2\).
This prefix ends at an equal-run boundary, including when some
atoms are exactly at the saturation threshold. Rounding the square
root in \(H_S\) costs at most a factor 2, so that
\[
 \widehat X(n):=\min_S\max\{Q_S,h_Sn^2\}\in[X(n)/16,X(n)].
\]
If the tail mass is zero, all these tail moments and \(X\) vanish.
Taking the maximum with \(\Delta\) gives
\begin{equation}
  Y(n)/17\le\widehat Y(n)=\max\{\Delta,\widehat X(n)\}\le Y(n),
  \qquad Y(n)=\Delta+X(n).
  \label{eq:closed-tail-comparison}
\end{equation}

For \(n=4^k\), the condition \(\widehat D(n)\le\widehat Y(n)\)
means that some \(R\) satisfies, for every \(S\) and \(a=0,1,2\),
at least one of
\[
 c_{R,a}\le b_S2^{ak},
 \qquad c_{R,a}\le h_S2^{(a+4)k}.
\]
Each inequality is an upward interval of nonnegative integer
exponents. Its least exponent is \(E\). The first exponent in
the union of two such intervals is the smaller of the two
\(E\)'s, and their intersection over the requirements has first
exponent equal to the maximum. Minimizing over \(R\) gives
\eqref{eq:closed-rate}. At the full head cut the constant
omitted-mass term is zero, so a finite exponent always exists:
for every tail cut at least one budget is positive when \(e>0\).

A successful closed comparison enforces the baseline scales
even without an explicit lower search floor. The full tail cut
forces \(1/n\le\max(\Delta,w)\le e\).
For the other scale, choose \(S=\{i\in T:q_i\ge\delta\}\).
In the regular branch \(q(S)\le J<32\delta\), and
\(\Delta<\min_{i\in H}q_i\le\max_iq_i\le\max(J,\delta)\).
The remaining moment is at most \(\sqrt\delta J\). The \(1/n\)
term thus forces
\[
 1/n\le\max\{32\delta,1024\delta^3n^2\},
\]
and hence \(n\delta\ge1/32\).

By \eqref{eq:closed-head-comparison}--\eqref{eq:closed-tail-comparison},
passing the closed comparison implies \(D(n)\le8Y(n)\).
Conversely, \(D(n)\le Y(n)/17\) suffices to pass. The horizontal
comparison and the statistical bounds therefore identify
\eqref{eq:closed-rate} with the sample complexity up to universal
constants at the same radius. A fixed dilation covers the baseline
floor, and rounding a clock up to the next power of four costs
at most a factor four.

\subsection{Exact arithmetic and scope}

A finite input may specify distinct rational atom weights together
with their binary multiplicities. Sort the weights, and locate the
shortest head by one rational ceiling if it ends within a run.
Delete one largest atom for the head cost. The remaining head
and tail scans use run endpoints. The displayed sums are computed
using multiplicities and do not enumerate the labels in a run.

A downward rounding \(\lfloor x^{u/v}\rfloor_2\) for the fixed
rational powers used here is computed by comparing the numerator and
denominator of \(x^u\) with integer shifts. Similarly, the ceiling binary
logarithm of a positive rational ratio is its numerator-denominator
bit-length difference adjusted by one exact comparison. The
definition of \(E\) then uses integer ceiling division. Common
denominators have bit length bounded by the sum of the input
denominator bit lengths, and each fixed-power and shift step
preserves polynomial bit length. The exponent and the binary
output \(4^k\) have polynomial bit length. The number of head-tail
cut pairs is quadratic in the number of input runs. Thus
\eqref{eq:closed-rate} is evaluable in polynomial bit time without
a root computation, a scan over sample sizes, or expansion of binary
multiplicities. This bit bound concerns evaluation of the rate.
The statistical theorem bounds the number of observations.

For a countable alphabet, the next section gives an exact projection argument
and relates the direct profile to its finite evaluations.
